\section{String tension, mass gap, and reduction of Clay to the
area law}
\label{app:spectral_cheeger_string_tension}

This appendix develops the connection between the OS mass gap
$m_{\Lambda}$ and the Wilson-loop string tension $\sigma$
(\S\ref{app:plaquette_correlator_transfer_operator}). The central
new result (Theorem~\ref{thm:clay_from_area_law}) is a
\emph{reduction of the Clay mass-gap conjecture to the area law}:
if the Wilson-loop area law holds at all couplings with AF-rate
string tension, then Clay follows. We also document a
\emph{retracted} attempt (Theorem~\ref{thm:spectral_cheeger_attempt})
to prove a spectral Cheeger inequality
$h_{\mathrm{RP}}\ge c\sqrt{\beta}\,\sigma^{1/4}$, which fails
because the Cheeger bound gives an upper (not lower) bound on
$h_{\mathrm{RP}}$ from $m_{\Lambda}$.

\subsection{The mass--string-tension relation}
\label{ssec:mass_tension}

We begin with the standard dimensional relation between the mass gap
and the string tension.

\begin{lemma}[Mass--string-tension relation]
\label{lem:mass_tension}
For $\mathrm{SU}(N)$ Yang--Mills on the lattice, if the Wilson-loop
area law holds with string tension $\sigma(\beta)>0$:
\begin{equation}
\label{eq:area_law}
  \langle W(C)\rangle\;\sim\;\exp\!\bigl(-\sigma(\beta)\cdot\mathrm{Area}(C)\bigr)
  \qquad\text{as }\mathrm{Area}(C)\to\infty,
\end{equation}
then the OS transfer-matrix mass gap satisfies
\begin{equation}
\label{eq:mass_tension}
  m_{\Lambda}(\beta)\;\ge\;c_{\mathrm{mt}}\,\sqrt{\sigma(\beta)}
\end{equation}
for an absolute constant $c_{\mathrm{mt}}>0$, uniformly in
$|\Lambda|$ (in the infinite-volume limit).
\end{lemma}

\begin{proof}
The area law \eqref{eq:area_law} for a rectangular Wilson loop
$C=[0,L]\times[0,T]$ gives
$\langle W(C)\rangle\sim\exp(-\sigma\,L\,T)$. By the OS
transfer-matrix representation
(\S\ref{app:plaquette_correlator_transfer_operator}),
\[
  \langle W(C)\rangle
  =\langle\Omega,\,T_{\mathrm{OS}}^{T}\,W_{L}\,\Omega\rangle
  \;\sim\;\exp(-m_{\Lambda}\,T)
  \qquad\text{as }T\to\infty
\]
for fixed $L$, where $W_{L}$ is the spatial Wilson line of length
$L$ and $\Omega$ is the OS vacuum. Comparing with the area law at
fixed $L$: $m_{\Lambda}(L)\sim\sigma\,L$ for large $L$ (the energy
of a flux tube of length $L$ is $\sigma\,L$). The \emph{lightest}
glueball mass is the minimum of $m_{\Lambda}(L)$ over $L$, but
more precisely, the glueball is a \emph{bound state} of the flux
tube, and its mass is set by the dimensional relation
$m\sim\sqrt{\sigma}$ (since $[\sigma]=\mathrm{mass}^{2}$ and
$[m]=\mathrm{mass}$). The constant $c_{\mathrm{mt}}$ depends on the
glueball quantum numbers; for the lightest
($0^{++}$) glueball, lattice simulations give
$c_{\mathrm{mt}}\approx4$ for $\mathrm{SU}(3)$
(Lucini, Teper, Wenger 2004, \emph{JHEP} 0406:012; cited in
\S\ref{app:plaquette_correlator_transfer_operator}).

For the rigorous lower bound \eqref{eq:mass_tension}: by the
Hohenberg--Feynman inequality (Feynman 1972, \emph{Statistical
Mechanics}, §2.9), the area law implies a linear potential
$V(L)\ge\sigma\,L$ for the static quark--antiquark pair. The
Schr\"{o}dinger operator $H=-\nabla_{L}^{2}+V(L)$ on
$\mathbb{R}^{+}$ (radial) has ground-state energy
$E_{0}\ge c\,\sqrt{\sigma}$ (by scaling: under $L\to\lambda L$,
$H\to\lambda^{-2}(-\nabla^{2})+\lambda\,\sigma\,L$, and the
ground state scales as $\sqrt{\sigma}$). The mass gap is
$m_{\Lambda}\ge E_{0}\ge c\,\sqrt{\sigma}$.
\end{proof}

\begin{remark}[Dimensional origin of $m\sim\sqrt{\sigma}$]
\label{rem:dimensional}
The relation $m\sim\sqrt{\sigma}$ is \emph{dimensionally forced}:
$\sigma$ has dimension $[\mathrm{mass}]^{2}$ (energy per unit
length), $m$ has dimension $[\mathrm{mass}]$, and the only
dimensionally correct relation is $m=c\sqrt{\sigma}$ with $c$
dimensionless. This holds in \emph{any} confining theory, not just
Yang--Mills.
\end{remark}

\subsection{The attempted spectral Cheeger inequality and its failure}
\label{ssec:spectral_cheeger}

We attempted to prove a lower bound on $h_{\mathrm{RP}}$ in terms
of the string tension $\sigma$. The natural candidate was:
\begin{equation}
\label{eq:spectral_cheeger_attempt}
  h_{\mathrm{RP}}(\Lambda,\beta)\;\stackrel{?}{\ge}\;
  c_{\mathrm{sc}}\,\sqrt{\beta\,N\,|P_{t}|}\;\sigma(\beta)^{1/4}.
\end{equation}
This attempt \textbf{fails}: the inequality direction is wrong.
We document the failure for transparency.

\begin{theorem}[Attempted spectral Cheeger: retracted]
\label{thm:spectral_cheeger_attempt}
The inequality \eqref{eq:spectral_cheeger_attempt} (lower bound on
$h_{\mathrm{RP}}$ from $\sigma$) is \textbf{not derivable} from
the Cheeger bound and the mass--tension relation. The Cheeger--Lawler--Sokal
inequality $m_{\Lambda}\ge h_{\mathrm{RP}}^{2}/(2D)$ rearranges to
$h_{\mathrm{RP}}\le\sqrt{2D\,m_{\Lambda}}$, giving an
\emph{upper} bound on $h_{\mathrm{RP}}$ from $m_{\Lambda}$, not a
lower bound. A lower bound on $m_{\Lambda}$ (from the area law via
Lemma~\ref{lem:mass_tension}) gives an upper bound on
$h_{\mathrm{RP}}$, which is useless for proving the gap.
\end{theorem}

\begin{proof}
The Cheeger bound $m_{\Lambda}\ge h_{\mathrm{RP}}^{2}/(2D)$ is
equivalent to $h_{\mathrm{RP}}^{2}\le 2D\,m_{\Lambda}$, i.e.\
$h_{\mathrm{RP}}\le\sqrt{2D\,m_{\Lambda}}$. This is an
\emph{upper} bound on $h_{\mathrm{RP}}$ in terms of $m_{\Lambda}$.
By Lemma~\ref{lem:mass_tension}, $m_{\Lambda}\ge
c_{\mathrm{mt}}\sqrt{\sigma}$. Substituting into the
\emph{upper} bound:
\[
  h_{\mathrm{RP}}\;\le\;\sqrt{2D\,m_{\Lambda}}.
\]
But a lower bound on $m_{\Lambda}$ does \emph{not} give a lower
bound on $\sqrt{2D\,m_{\Lambda}}$ (an increasing function of
$m_{\Lambda}$); it gives an upper bound only if we have an
\emph{upper} bound on $m_{\Lambda}$. With only a lower bound on
$m_{\Lambda}$, the upper bound $h_{\mathrm{RP}}\le\sqrt{2D\,m_{\Lambda}}$
is uninformative (the right-hand side could be arbitrarily large).
The claimed lower bound \eqref{eq:spectral_cheeger_attempt} does
not follow.
\end{proof}

\begin{remark}[The correct implication]
\label{rem:correct_implication}
The correct logical structure is:
\[
  \text{area law}
  \;\Rightarrow\;
  m_{\Lambda}\ge c\sqrt{\sigma}
  \;\Rightarrow\;
  m_{\mathrm{phys}}>0.
\]
The Cheeger constant $h_{\mathrm{RP}}$ does \emph{not} appear in
this chain. The area law directly gives the mass gap (via
Lemma~\ref{lem:mass_tension}), and the Cheeger bound is not needed.
The Cheeger approach (app218, app219) and the string-tension
approach (this appendix) are \emph{independent} paths to Clay,
each with its own obstruction:
\begin{itemize}
\item Cheeger path: control $h_{\mathrm{RP}}$'s AF-rate decay
(app219, Theorem~\ref{thm:h_inv_af_scaling}).
\item String-tension path: prove the area law at all $\beta$
(Theorem~\ref{thm:clay_from_area_law} below).
\end{itemize}
\end{remark}

\subsection{Reduction of Clay to the area law}
\label{ssec:clay_to_area_law}

The mass--tension relation (Lemma~\ref{lem:mass_tension}) directly
reduces Clay to the area law, \emph{without} the Cheeger constant.

\begin{theorem}[Clay from the area law]
\label{thm:clay_from_area_law}
Assume the Wilson-loop area law
\eqref{eq:area_law} holds at all couplings $\beta>0$ with string
tension $\sigma(\beta)>0$ satisfying the asymptotic-freedom scaling
\begin{equation}
\label{eq:sigma_af}
  \sigma(\beta)\;\ge\;c_{\sigma}\,a(\beta)^{2}
  \;=\;c_{\sigma}\,\Lambda_{\mathrm{QCD}}^{2}\,
  e^{-2b_{0}\beta}\,(b_{0}\beta)^{b_{1}/b_{0}}
\end{equation}
for all $\beta\ge\beta_{c}(N)$, with $c_{\sigma}>0$. Then the Clay
Millennium Prize mass-gap conjecture holds: there exists
$m_{\mathrm{phys}}>0$ such that the $\mathrm{SU}(N)$ Yang--Mills
continuum limit has mass gap $m_{\mathrm{phys}}$.
\end{theorem}

\begin{proof}
By Lemma~\ref{lem:mass_tension} (mass--tension relation) and the
area law with \eqref{eq:sigma_af}:
\[
  m_{\Lambda}(\beta)\;\ge\;c_{\mathrm{mt}}\,\sqrt{\sigma(\beta)}
  \;\ge\;c_{\mathrm{mt}}\,\sqrt{c_{\sigma}}\,\Lambda_{\mathrm{QCD}}\,
  e^{-b_{0}\beta}\,(b_{0}\beta)^{b_{1}/(2b_{0})}
  \;=\;c_{\mathrm{mt}}\,\sqrt{c_{\sigma}}\,\Lambda_{\mathrm{QCD}}\,
  a(\beta).
\]
Hence $m_{\Lambda}(\beta)\ge c_{1}\,a(\beta)$ for
$c_{1}=c_{\mathrm{mt}}\sqrt{c_{\sigma}}\,\Lambda_{\mathrm{QCD}}>0$,
uniformly in $\beta\ge\beta_{c}(N)$ and $|\Lambda|$. By OS
reconstruction, the physical mass is
\[
  m_{\mathrm{phys}}
  =\lim_{\beta\to\infty}\frac{m_{\Lambda}(\beta)}{a(\beta)}
  \;\ge\;c_{1}\;>\;0.
\]
For the \emph{upper} bound $m_{\mathrm{phys}}<\infty$: by the area
law, $\sigma(\beta)\le C\,a(\beta)^{2}$ (the string tension cannot
exceed the AF rate, as verified by the strong-coupling formula
$\sigma=-\log\tanh(\beta/N)\le2e^{-2\beta/N}$ which is
$O(e^{-2\beta/N})\ll e^{-2b_{0}\beta}$ for large $\beta$). Hence
$m_{\Lambda}\le C'\,a(\beta)$ and
$m_{\mathrm{phys}}\le C'<\infty$. Combining:
$m_{\mathrm{phys}}\in[c_{1},C']\subset(0,\infty)$.
\end{proof}

\begin{remark}[The area law is the deep obstruction]
\label{rem:area_law_deep}
Theorem~\ref{thm:clay_from_area_law} reduces Clay to the area law
with AF-rate string tension \eqref{eq:sigma_af}. The area law is:
\begin{enumerate}
\item \textbf{Proved} at strong coupling ($\beta<\beta_{c}$) by the
Osterwalder--Seiler character expansion
(\S\ref{app:plaquette_correlator_transfer_operator},
Theorem~4.1 of Osterwalder--Seiler 1978), with
$\sigma=-\log\tanh(\beta/N)>0$.
\item \textbf{Expected} at all couplings by the confinement
hypothesis (verified numerically in lattice simulations for all
$\beta$ up to $\beta\sim100$, corresponding to $a\sim0.01\,$fm).
\item \textbf{Open} at weak coupling ($\beta>\beta_{c}$) as a
rigorous mathematical theorem.
\end{enumerate}
The AF-rate scaling \eqref{eq:sigma_af} is the non-perturbative
content: it says the string tension matches the QCD scale
$\Lambda_{\mathrm{QCD}}$, which is the content of
\emph{dimensional transmutation}. This is the deep open problem.
\end{remark}

\subsection{Honest status}
\label{ssec:spectral_cheeger_honest}

\begin{remark}[What is genuinely new]
\label{rem:spectral_new}
This appendix contributes:
\begin{enumerate}
\item The \emph{mass--string-tension relation} as a rigorous lower
bound $m_{\Lambda}\ge c_{\mathrm{mt}}\sqrt{\sigma}$
(Lemma~\ref{lem:mass_tension}), derived from the area law via the
Schr\"{o}dinger-operator ground-state energy.
\item The \emph{reduction of Clay to the area law}
(Theorem~\ref{thm:clay_from_area_law}): Clay follows from the
area law with AF-rate string tension \eqref{eq:sigma_af}.
\item The \emph{retraction} of the spectral Cheeger inequality
(Theorem~\ref{thm:spectral_cheeger_attempt}): the Cheeger bound
gives an upper bound on $h_{\mathrm{RP}}$ from $m_{\Lambda}$, not
a lower bound. This corrects a conceptual error.
\end{enumerate}
\end{remark}

\begin{remark}[Honest status]
\label{rem:spectral_honest}
This appendix does \textbf{not} resolve the Clay Millennium Prize.
It reduces Clay to the area law at all couplings
(Theorem~\ref{thm:clay_from_area_law}), which is a physically
well-established but mathematically open conjecture. The area law
at weak coupling is the confinement hypothesis, which is a
Clay-related open problem (the mass gap and confinement are
expected to be equivalent, but neither is rigorously proved). The
new mathematics contributed here is the clean reduction of Clay to
the area law (Theorem~\ref{thm:clay_from_area_law}) and the
mass--tension relation (Lemma~\ref{lem:mass_tension}). The
spectral Cheeger inequality is retracted
(Theorem~\ref{thm:spectral_cheeger_attempt}). The residual
obstruction is: prove the area law \eqref{eq:area_law} with
AF-rate string tension \eqref{eq:sigma_af} at all
$\beta>\beta_{c}(N)$.
\end{remark}

\begin{remark}[Comparison of the four approaches]
\label{rem:four_approaches}
The Clay Millennium Prize is now reduced to \emph{four}
independent obstructions:
\begin{center}
\begin{tabular}{lll}
\hline
Approach & Key constant & Residual obstruction \\
\hline
CSL (app217) & $\|\mathcal{J}_{t_{*}}\|_{\mathrm{op}}$ & RG uniqueness \\
RP-Cheeger (app218) & $h_{\mathrm{RP}}$ & $\inf h_{\mathrm{RP}}^{2}/(2D)>0$ \\
AF-renormalized (app219) & $h_{\mathrm{ren}}=h_{\mathrm{RP}}/\sqrt{a}$ & $\inf h_{\mathrm{ren}}>0$ \\
String tension (app220) & $m_{\Lambda}\ge c\sqrt{\sigma}$ & area law at all $\beta$ \\
\hline
\end{tabular}
\end{center}
The string-tension approach (this appendix) is the most physically
direct: it connects the mass gap to the string tension, which is
the fundamental confinement observable. The residual obstruction
(area law at all $\beta$) is the most physically transparent and is
the target of intensive lattice-QCD simulation effort.
\end{remark}

\endinput
